\clearpage
\fancyhead[L]{\small KOMET.}
\SourceLeaf{447}
\begin{center}
{\large KOMET.}

\noindent\rule{\linewidth}{0.4pt}

Göttingische gelehrte Anzeigen. Nr. 26. S. 249..251. 15. Februar 1823.

\noindent\rule{\linewidth}{0.4pt}
\end{center}

Ein Schreiben von Herrn \textsc{Rümker} aus Paramatta in New South Wales an Hofrat \textsc{Gauss} vom August 1822 enthält neben vielen anderen dort angestellten astronomischen Beobachtungen, die erst an anderer Stelle veröffentlicht werden sollen, auch die Beobachtungen des Kometen von \textsc{Encke}, die Herr \textsc{Rümker} im Juni des vergangenen Jahres angestellt hat. Bekanntlich war die Hoffnung sehr gering, den Kometen bei seiner Wiederkehr zur Sonne im Jahr 1822 in Europa beobachten zu können. Tatsächlich blieben alle Bemühungen europäischer Astronomen, ihn zu sehen, völlig erfolglos. Dagegen rechnete man damit, dass er den Beobachtern auf der südlichen Halbkugel nicht entgehen würde. Da aber bereits bekannt war, dass die Nachforschungen der englischen Astronomen am Kap der Guten Hoffnung \emph{nicht gelungen} waren, wartete man umso sehnlicher auf Nachrichten aus Neuholland. Umso mehr freuen wir uns, in unseren Blättern die Bestätigung mitteilen zu können, die unser geschickter Landsmann gewonnen hat. In Nr. 83 des Jahrgangs 1819 [S.~420 dieses Bandes.] haben wir zuerst über die höchst bemerkenswerte Entdeckung von Herrn \textsc{Encke} berichtet. Es folgen \emph{sämtliche} Beobachtungen von Herrn \textsc{Rümker}. Nach dem 23. Juni störte das Mondlicht, und nach dem Vollmond war der Komet für weitere Beobachtungen zu lichtschwach.

\Needspace{23\baselineskip}
\begin{center}
Beobachtungen des Kometen von \textsc{Encke} in \emph{Paramatta}.
\end{center}

\begin{center}
\small\setlength{\tabcolsep}{3pt}
\begin{tabular}{l|ccc|ccc|ccc}
\multicolumn{1}{c|}{1822} & \multicolumn{3}{c|}{Sternzeit} & \multicolumn{3}{c|}{Mittlere Rektasz.} & \multicolumn{3}{c}{Mittlere Deklination} \\
\hline
Juni 2 & $10^{\mathrm{u}}$ & $39^{\mathrm{m}}$ & $25^{\mathrm{s}}$ & $92^{\circ}$ & $43'$ & $51''3$ & $17^{\circ}$ & $39'$ & $46''3$ Nordl. \\
3 & 11 & --- & --- & 93 & 46 & 20.7 & 16 & 53 & 7.5 \\
4 & 11 & 3 & 0 & 94 & 46 & 0.0 & 16 & 4 & 36.7 \\
6 & 11 & 7 & 38 & 96 & 42 & 11.6 & 14 & 22 & 42.0 \\
7 & 11 & 3 & 10 & 97 & 38 & 15 & 13 & 26 & 5 \\
8 & 11 & 17 & 25 & 98 & 33 & 47.7 & 12 & 31 & 18.6 \\
10 & 11 & 20 & 0 & 100 & 24 & 43.8 & 10 & 29 & 49.5 \\
11 & 11 & 24 & 39 & 101 & 19 & 44.5 & 9 & 26 & 4.6 \\
12 & 11 & 40 & 0 & 102 & 17 & 52 & 8 & 18 & 30 \\
13 & 11 & 42 & 4 & 103 & 15 & 2 & 7 & 6 & 30 \\
14 & 11 & 55 & 0 & 104 & 15 & 40 & 5 & 52 & 27 \\
15 & 11 & 40 & 48 & 105 & 17 & 0.5 & 4 & 33 & 40 \\
19 & 12 & 13 & 38 & 109 & 54 & 36.4 & 1 & 29 & 43.7 Südl. \\
20 & 12 & 16 & 53 & 111 & 14 & 26.9 & 3 & 14 & 29.1 \\
22 & 13 & 18 & 46 & 114 & 12 & 20.5 & 7 & 8 & --- \\
23 & 12 & 53 & 55 & 115 & 47 & 41.7 & 9 & 9 & 48.4 \\
\end{tabular}
\end{center}

Es wird nicht gesagt, ob der Zusatz \emph{mittlere} bei Rektaszension und Deklination ein Mittel aus mehreren Beobachtungen meint oder eine Bereinigung um die Nutation. Zum Vergleich wurden Sterne aus dem Katalog von \textsc{Piazzi} und aus der \emph{Histoire Céleste} verwendet. Herr \textsc{Rümker} verspricht, Näheres bei nächster Gelegenheit nachzuliefern.

Aus den beiden Sonnenwenden im Dezember 1821 und im Juni 1822 erhält Herr \textsc{Rümker} für die Sternwarte in Paramatta folgende Polhöhe

\[
33^{\circ}\,48'\,41''97
\]

\Needspace{16\baselineskip}
Aus einer Bedeckung von $\alpha$ Scorpii am 10. April 1822, bei der

% 
\SourceLeaf{448}
\begin{center}
\small
\begin{tabular}{l l}
der Eintritt \mbox{.\,.\,.\,.\,.\,.\,.\,.\,.\,.\,.} & $18^{\mathrm{u}}\;35^{\mathrm{m}}\;47^{\mathrm{s}}4$ Mittlere Zeit \\
der Austritt \mbox{.\,.\,.\,.\,.\,.\,.\,.\,.\,.\,.\,.} & $19\quad14\quad27.9$
\end{tabular}
\end{center}

\noindent beobachtet wurden, bestimmt Herr \textsc{Rümker} mit den dem Nautical Almanac entnommenen Positionen des Mondes die Länge von Paris
\begin{center}
\begin{tabular}{l l}
10 St. $3^{\mathrm{m}}\;56^{\mathrm{s}}3$ & aus dem Eintritt \\
$10\quad4\quad7.3$ & aus dem Austritt.
\end{tabular}
\end{center}
